\documentclass[a4paper,10pt]{letter}

\usepackage[T1]{fontenc}
\usepackage[utf8]{inputenc}
\usepackage{lmodern}
\usepackage[hidelinks]{hyperref}
\usepackage{microtype}
\usepackage[margin=1in]{geometry}

\signature{Omar Iskandarani\\
Independent Researcher, Groningen\\
The Netherlands\\
ORCID: 0009-0006-1686-3961\\
Email: \href{mailto:info@omariskandarani.com}{info@omariskandarani.com}}

\address{Omar Iskandarani\\
Vinkenstraat 86A\\
9713 TK Groningen\\
The Netherlands}

\date{\today}

\begin{document}
    
    \begin{letter}{Editors\\
    \textit{European Journal of Physics}}
        
        \opening{Dear Editors,}
        
        Please consider the enclosed manuscript entitled
        \textbf{“Hydrogen Ground-State Hyperfine Splitting as a Canonical Activation Gap: From the Fermi Contact Interaction to the Schottky Heat-Capacity Anomaly”}
        for publication in \textit{European Journal of Physics}.
        
        This manuscript presents a self-contained derivation chain connecting two standard topics that are usually taught separately: the hydrogen $1s$ hyperfine splitting derived from the Fermi contact interaction, and the Schottky heat-capacity anomaly of the resulting two-level system with degeneracies $(g_0,g_1)=(1,3)$. The paper is fully orthodox in its physics and does not claim new fundamental results. Its contribution is pedagogical and synthetic: it provides a logically explicit bridge from atomic magnetic structure to low-temperature thermodynamics in a single worked derivation.
        
        In particular, the manuscript derives the leading-order hyperfine gap in closed form, relates it to the experimentally observed 21\,cm transition, and then propagates that same gap into the exact canonical heat capacity, including the low-temperature asymptotic regime and an explicit derivation of the peak condition for the degeneracy-asymmetric case. The goal is to make transparent how a spectroscopic energy scale becomes a thermodynamic activation scale, and to do so in a form suitable for advanced undergraduate and graduate readers.
        
        We believe the manuscript is well suited to \textit{European Journal of Physics}, given the journal’s emphasis on clear, rigorous, and pedagogically valuable presentations of established physics. The paper is intended as a compact teaching-oriented research note in that spirit: not a claim of new physics, but a worked synthesis that may be useful in courses and self-study at the interface of atomic physics, statistical mechanics, and thermal physics.
        
        This manuscript is original, has not been published previously, and is not under consideration for publication elsewhere.
        
        Thank you for your time and consideration.
        
        \closing{Sincerely,}
    
    \end{letter}

\end{document}